% Generated draft: requires mathematical and source review.
\documentclass[11pt]{article}
\usepackage[margin=25mm]{geometry}
\usepackage{fontspec}
\setmainfont{Latin Modern Roman}
\usepackage{amsmath,amssymb,mathtools,unicode-math}
\setmathfont{Latin Modern Math}
\usepackage{enumitem,needspace,xurl,hyperref}
\setlength{\parindent}{0pt}
\setlength{\parskip}{5pt}
\setlength{\emergencystretch}{3em}
\newcommand{\N}{\mathbb{N}}
\newcommand{\Z}{\mathbb{Z}}
\newcommand{\Q}{\mathbb{Q}}
\newcommand{\R}{\mathbb{R}}
\newcommand{\C}{\mathbb{C}}
\newcommand{\F}{\mathbb{F}}
\newcommand{\A}{\mathbb{A}}
\newcommand{\PP}{\mathbb{P}}
\newcommand{\E}{\mathbb{E}}
\newcommand{\eps}{\varepsilon}
\newcommand{\dd}{\,\mathrm{d}}
\begin{document}
\section*{UnsolvedMath ID: 30001749}
\textbf{Source problem code:} catalogue entry
\subsection{Definitions and mathematical statement}
Let \(d,m\) be positive integers and let \(\rho=(\rho_1,\ldots,\rho_m)\in(0,1)^m\) satisfy
\[
\sum_{j=1}^{m}\rho_j^d<1.
\]
A self-similar realization with contraction vector \(\rho\) is a nonempty compact set \(E\subseteq\R^d\) satisfying
\[
E=\bigcup_{j=1}^{m}S_j(E),\qquad
\|S_j(x)-S_j(y)\|=\rho_j\|x-y\|\quad(x,y\in\R^d),
\]
where the \(S_j\) are similarities. We use the usual strong separation interpretation of \emph{dust-like}: \(S_i(E)\cap S_j(E)=\varnothing\) whenever \(i\ne j\). This interpretation is a drafting convention because the supplied excerpt does not explicitly define ``dust-like.'' Write \(E=D(\rho)\) for a chosen such realization; the notation does not specify its defining similarities.

Let \(\tau=(\tau_1,\ldots,\tau_n)\in(0,1)^n\) be another finite contraction vector and let \(F=D(\tau)\) be an associated dust-like self-similar realization. Equip both sets with their induced Euclidean metrics, denoted by \(d_E\) and \(d_F\). The supplied statement does not explicitly specify the ambient dimension of \(F\), or whether an inequality analogous to \(\sum_j\rho_j^d<1\) is imposed on \(\tau\); these admissibility details remain unresolved.

The sets \(E\) and \(F\) are \emph{Lipschitz equivalent} if there exist a bijection \(f:E\to F\) and a constant \(L\geq1\) such that
\[
L^{-1}d_E(x,y)\leq d_F(f(x),f(y))\leq Ld_E(x,y)
\qquad\text{for all }x,y\in E.
\]
\textbf{Problem.} Characterize nontrivial necessary and sufficient conditions on the contraction vectors \(\rho\) and \(\tau\) for their associated dust-like self-similar sets \(D(\rho)\) and \(D(\tau)\) to be Lipschitz equivalent. The wording ``any'' in the supplied statement suggests that the criterion should apply to arbitrary admissible realizations. The excerpt does not explicitly settle whether this means a criterion uniform over all choices of defining similarities or a question about particular chosen realizations. This quantifier qualification must be resolved before treating the formulation as definitive. The requested word ``nontrivial'' is also not given a formal meaning.
\subsection{Short English statement}
Given two contraction vectors and associated dust-like self-similar sets, seek substantive conditions on the vectors that are equivalent to the existence of a bi-Lipschitz bijection between the sets. The precise admissibility hypotheses for the second vector, the meaning of dust-like, and the quantification over realizations require confirmation from the source.
\subsection{Sources}
\begin{itemize}
\item Local problem registry: \texttt{data/unsolvedmath/problems.json}, with \texttt{catalog/problems.json} as the supplementary registry; supplied permanent UnsolvedMath ID \(30001749\), exact source problem code \texttt{catalogue entry}, title \emph{Lipschitz Equivalence of Dust-Like Self-Similar Sets}. This draft uses the supplied complete record; the registry files were not independently inspected.
\item Supplied report citation: \emph{Operator Algebras and Representation Theory: Frames, Wavelets and Fractals} (2011), Oberwolfach Reports, pp.~901--978, DOI \texttt{10.4171/owr/2011/17}. Source URL and supplied literature reference labeled ``original OWR report'': \url{https://doi.org/10.4171/owr/2011/17}. The record attributes the problem to workshop open-problem discussions and gives proposed year 2011; no individual proposer is supplied.
\item Supplied literature assessment, checked \texttt{2026-08-21}: the record labels the problem open and says that the extracted notation and hypotheses are truncated and that no resolution of the general equivalence question is identifiable there. This is supplied dated curation, not an independently verified literature assessment.
\item The supplied statement status is \texttt{corrected\_verified}; its verification note says that the official report was checked and the contraction-vector hypotheses and meaning of the associated dust-like sets were restored. That supplied verification claim has not been independently checked for this draft and remains in tension with the truncation qualification in the supplied literature assessment.
\end{itemize}
\end{document}
